\subsection{Solving Equations with Rational Powers} 
We can generalize the way that we solve power equations to solve equations involving rational exponents, using the following facts.
\begin{fact}
If $X$ is an expression that represents a positive real number,
$$\left(X^{p/q}\right)^{q/p}=X$$
for all values of the variables in $X$.
\end{fact}
We know this is true because, when all the expressions involved are defined, and we only take even powers of positive numbers,
\begin{lpic}{}{1}\end{lpic}
Note that here, we're restricting ourselves to \makeblank numbers. This works for all real numbers when $p$ and $q$ are both \makeblank[0.75in], but we need to worry about signs if $p$ or $q$ is \makeblank[0.75in].
\begin{fact}
If $X$ and $Y$ represent positive real numbers, and $\frac{p}{q}$ is a rational number written in lowest terms the equations
$$X=Y
\quad\mbox{and}\quad
X^{p/q}=Y^{p/q}$$
are equivalent.
\end{fact}
\begin{fact}
If $p$ and $q$ are both odd numbers, the above facts hold even if $X$ and $Y$ are negative.
\end{fact}
\begin{example}
Use the facts above to find all positive real solutions to each equation below.
\begin{lpicvsplit}{3}{6}
\foreach \x/\eqn in {0/{x^{2/3}=9},1/{3x^{3/4}=81},2/{(x+1)^{3/5}=8}}
	\regtext{\x*\value{fcols}+\x+0.5*\value{fcols}}{-1}{$\eqn$};
\end{lpicvsplit}
\end{example}